% Paper2 figure-led manuscript draft — 2026-05-11
% Evidence boundary: only the 2026-05-11 Remote107 selective-KV claim-lock package is claim-bearing.
% Superseded candidate TSVs and Gemini-era figures are audit-only unless regenerated from that package.

\section*{Selective Analog KV Cache for Session-Lifetime Language-Model Memory}

\begin{abstract}
Decoder-only language models store Key/Value tensors that are written during prefill or decoding, read repeatedly during one session, and discarded. This lifetime is different from both permanent model weights and short-lived activation buffers, making KV cache a distinct memory workload. We study whether analog memory should be used for the whole KV cache or only for selected layers. Under the locked Remote107 protocol for Pythia-410M on WikiText-2 raw-v1, last1/last2 terminal-layer scopes remain viable while all-layer analog KV degrades sharply. The result supports a selective analog cache tier and does not claim energy, endurance, refresh, long-context, or larger-model closure.
\end{abstract}

\paragraph{Central claim.}
A decoder KV cache is not a smaller version of model-weight storage. Each entry is created during prefill or decoding, read repeatedly during one session, and discarded. This access pattern makes KV cache a candidate for session-lifetime analog memory tiers. The present claim is deliberately narrow: under the locked Remote107 protocol, last1/last2 terminal-layer analog KV remains viable, whereas all-layer analog KV degrades sharply and should be used only as a negative control.

\section{Introduction}

Long-context inference turns KV cache into a dominant memory object. The cache grows linearly with context and batch size, and it is accessed differently from both weights and activations. Weights are persistent and repeatedly reused across requests; activations are transient within a forward pass. KV entries occupy the middle: they are generated once, kept for the session, read at every following token, and then released.

This access pattern suggests a different analog-memory question. The goal is not a fully analog language model or a replacement for digital compute. The question is narrower: can selected KV tensors be placed in a session-lifetime analog cache without corrupting perplexity, and which layer scope is safe enough to optimize further?

\section{Memory model}

Autoregressive inference repeatedly reads the accumulated Key/Value tensors. For a decoder with $L$ layers, $H$ heads, per-head dimension $D_h$, context length $S$, and element size $P$, the KV footprint is
\begin{equation}
M_{\mathrm{KV}} = 2 L H D_h S P.
\end{equation}
This footprint grows linearly with context and batch size. The write path is different from weights: each KV entry is created during prefill or decoding, read repeatedly during the session, and then released. The useful lifetime is therefore seconds to minutes, not years.

That distinction changes the hardware question. Long-retention nonvolatile memories may be over-specified for KV cache, while SRAM is area-expensive at large contexts. Organic electrochemical memories have limited endurance and session-scale retention; those limits are constraints for weights, but they are a plausible fit for a write-once-read-many cache if the cache is selective and the protocol reports time-dependent cost explicitly.

\begin{figure}[t]
\centering
\includegraphics[width=\linewidth]{figures/fig_paper2_overview_selective_kv.pdf}
\caption{\textbf{KV cache is a better physical match than permanent model weights.} A decoder KV cache is written once, read repeatedly, and released at session end; that lifetime is closer to a session-scale analog memory tier than to archival weight storage. The proposed route does not claim a fully analog language model: only selected terminal-layer K/V tensors are mapped to an analog cache while softmax, control flow, and most transformer layers remain digital.}
\label{fig:paper2_overview}
\end{figure}

\section{Method: selective analog KV}

The first design choice is not device type or bit width; it is layer scope. We compare four scopes: last1, last2, last4, and all24. If terminal layers tolerate analog KV but all layers do not, the correct architecture is a selective cache tier rather than a whole-model analog cache.

The numerical rows below are taken from the 2026-05-11 Remote107 selective-KV claim-lock summary TSV in the manuscript source-data package. They intentionally do not use the older Fresh-D2D candidate index. Values below the digital reference should not be read as an analog-over-digital accuracy claim; the locked claim is layer-scope ordering under this protocol.

\begin{table}[t]
\centering
\caption{\textbf{Remote107 claim-lock evidence supports terminal-layer selectivity and rejects all-layer analog KV as the main route.} The values come from the 2026-05-11 Remote107 selective-KV claim-lock summary TSV: Pythia-410M on WikiText-2 raw-v1 test with context length 512, non-overlapping stride 512, and batch size 1. Digital PPL is 23.30. Each analog row has five seeds and committed per-row JSON artifacts. The normalized ratio is a within-protocol reference only; values below 1x do not constitute an analog-over-digital baseline claim.}
\label{tab:paper2_remote107_claimlock}
\begin{tabular}{lccc}
\toprule
Scope & D2D & PPL & PPL / digital ref. \\
\midrule
last1 & 0.02 & $18.781\pm0.023$ & 0.806x \\
last1 & 0.05 & $18.963\pm0.035$ & 0.814x \\
last2 & 0.02 & $18.594\pm0.014$ & 0.798x \\
last2 & 0.05 & $19.018\pm0.025$ & 0.816x \\
last4 & 0.02 & $20.402\pm0.069$ & 0.876x \\
last4 & 0.05 & $21.535\pm0.051$ & 0.924x \\
all24 & 0.02 & $25.985\pm0.148$ & 1.115x \\
all24 & 0.05 & $59.055\pm1.610$ & 2.535x \\
\bottomrule
\end{tabular}
\end{table}

\begin{figure}[t]
\centering
\includegraphics[width=0.95\linewidth]{source_data/remote107_selective_kv_claim_lock_20260511/figures/fig_107_selective_kv_claim_lock.pdf}
\caption{\textbf{Selective terminal-layer analog KV survives the locked protocol; all-layer analog KV does not.} The figure is regenerated from the verified Remote107 source-data package containing 41 manifest rows and 41 original JSON artifacts. These rows are claim-bearing for layer-scope ordering, but they do not establish energy, endurance, refresh, or long-context deployment claims.}
\label{fig:paper2_remote107_claimlock}
\end{figure}

\section{Locked Remote107 evidence}

The locked packet answers one question cleanly: which layer scopes remain plausible under the corrected-noise Remote107 protocol? The answer is selective terminal-layer KV. At both evaluated D2D levels, last1 and last2 stay tightly clustered, and all-layer analog KV is a high-exposure negative control. Last4 is still usable as an intermediate stress point. All24 exceeds the digital reference at D2D=0.02 and degrades sharply at D2D=0.05.

This result should be framed as a design-rule result, not as a final product claim. The manuscript claim is: \emph{selective terminal-layer analog KV is the route worth optimizing; all-layer analog KV should be treated as a negative control}. It is not yet a claim of end-to-end energy savings, refresh-safe operation, or long-context deployment.

The newer 107-clean validation lane is consistent with that claim but should remain explicitly secondary. In the frozen-target 410M adaptive-schedule study, fixed, cosine, and layer-wise training reach 18.75, 18.31, and 18.36 PPL, while reverse layer-wise degrades to 21.30. At larger scales, fixed/cosine/layer-wise nearly tie (2.8B: 12.68/12.68/12.70; 6.9B: 11.40/11.38/11.43), suggesting that schedule sensitivity weakens with scale. The clearest schedule-specific gain appears when the analog scope expands beyond last1: on Pythia-2.8B, \texttt{\detokenize{last2+cosine}} reaches 12.56 PPL, improving over \texttt{\detokenize{last2+fixed}} at 13.78 and slightly outperforming \texttt{\detokenize{last1+fixed}} at 12.68. These rows sharpen the engineering route, but they should stay out of the main claim table until their source-data package is frozen and audited to the same standard as the 2026-05-11 410M packet.

\section{Evidence boundary}

Local cache-path probes, retention/refresh scans, offline KV reconstruction sweeps, and the newer 107-clean scale-up or Qwen3-VL pushes are useful because they test plumbing and identify likely failure axes. They should stay out of the main claim table unless rerun or packaged under a locked protocol with dataset, command, commit, seed, checkpoint hash, and artifact completeness comparable to Remote107. The Gemini-era plots \texttt{\detokenize{fig_107_selective_kv}}, \texttt{\detokenize{fig_107_noise_sweep}}, and \texttt{\detokenize{fig_107_training_steps}} remain candidate/audit visuals unless regenerated from the 2026-05-11 source-data package.

\begin{table}[t]
\centering
\caption{\textbf{Evidence boundary for the current Paper2 manuscript.} Only the Remote107 claim-lock package is used for the main layer-scope claim; all other rows are retained to define what must not be pooled into that claim.}
\label{tab:paper2_evidence_boundary}
\begin{tabular}{lll}
\toprule
Evidence source & Current role & Manuscript use \\
\midrule
Remote107 2026-05-11 claim-lock package & Claim-bearing & Main layer-scope result \\
107-clean scale-up/Qwen3-VL updates & Active validation & Cite only after claim-lock packaging \\
107-clean adaptive schedules / R3 figures & Active validation & Use only with explicit provisional boundary \\
2026-05-10 Fresh-D2D TSV & Candidate/audit & Do not use as current claim table \\
Gemini-era 107 figures & Candidate/audit & Regenerate before claim use \\
Local cache-path probe & Engineering/debug & Methods or supplementary only \\
Offline KV reconstruction & Engineering/debug & Mechanism diagnostic only \\
Retention/refresh scan & Engineering/debug & Future protocol target \\
Older 107-clean candidate lane & Superseded & Do not pool or cite as final \\
\bottomrule
\end{tabular}
\end{table}

\section{Discussion}

The main lesson is that analog KV cache should be designed as a selective memory hierarchy. The terminal layers are attractive because they reduce analog exposure while preserving the cache entries closest to the final token distribution. The all-layer result is valuable precisely because it degrades: it rules out an over-broad analogization strategy and turns Paper2 into an architecture-selection paper rather than a generic analog-memory paper.

The next locked experiments should add time and scale without weakening the current evidence boundary. The latest 107-clean pushes point to 2.8B/6.9B and Qwen3-VL validation lanes, but those lanes should enter the manuscript only after they have source-data packages and manifests comparable to the 410M Remote107 packet. A retention-aware claim must jointly report PPL, $\tau$, refresh interval, write budget, and energy. Until then, organic memory should be described as a physically motivated candidate for session-lifetime KV cache, not as a proven deployment substrate.

\section{Contributions and limits}

The current manuscript contributes three bounded points. First, it reframes KV cache as a session-lifetime WORM memory workload rather than permanent weight storage. Second, it identifies layer scope as the primary architecture variable for analog KV cache. Third, it provides a locked Remote107 source-data package showing that terminal-layer selectivity and all-layer rejection are reproducible under the stated Pythia-410M/WikiText-2 protocol.

The limits are equally important. The locked packet does not close retention, refresh, energy, endurance, long-context, or larger-model deployment. Those axes should be evaluated in a follow-on protocol that reports PPL together with $\tau$, refresh interval, write budget, and system cost.

\section{Conclusion}

The current Paper2 story can be short and strong. KV cache has a session-lifetime WORM access pattern; organic analog memory is physically plausible for that workload; and the locked Remote107 packet shows that terminal-layer selectivity is the viable route while all-layer analog KV is rejected. The next manuscript step is to upgrade retention/refresh from engineering probes to a locked protocol.
